% 2026-11-21 / v1.40
% !TEX root = ../nc-quadric-classification.tex
\section{noncubic の良い根に対する中央 AS4 被覆}
\label{sec:noncubic-central-hull}

\paragraph{範囲.}
以下の構成は、中央二重平方環を与える従来の Ore 被覆とは異なる被覆である。
同じ trace 零交換対合で選んだ Serre 根の軌道環について、切断元自体が中央となる
AS4 被覆を与える。既存の Ore ambient の中で切断元を中央化する主張ではない。
新規性は未確認である。

基礎体を標数零の代数閉体とし、\(a,b\in k\)、
\(\theta=1-ab\ne0\)、\(\delta^2=\theta\) と置く。
二つの非退化二元二次形式は座標・定数倍を選べば
\[
 f_1=uv,\qquad f_2=(u-av)(bu-v)
       =bu^2-(1+ab)uv+av^2
\]
と書ける。二根対の一致も \(a=b=0\) に含まれる。
\[
 B_0=k[u,v]\langle c,d\rangle/(c^2-f_1,d^2-f_2)
\]
で \(u,v\) は中央とする。良い交換自己同型は
\[
 \sigma(u)=(u-av)/\delta,\quad
 \sigma(v)=(bu-v)/\delta,\quad
 \sigma(c)=d/\delta,\quad \sigma(d)=-\delta c
\]
であり、\(\sigma^2=\iota\)、\(\iota(u,v,c,d)=(u,v,-c,-d)\) である。
以下では \(x*y=x\sigma^{\deg x}(y)\) の Zhang 規約を使う。

\begin{theorem}\label{thm:noncubic-central-hull}
次の六関係環を定める。
\[
 A_0=T(u,v,c,d)\Big/
 \left(
 \begin{gathered}
 [u,c],[u,d],[v,c],[v,d],\quad c^2-\frac{uv+vu}{2},\\
 d^2-bu^2+\frac{1+ab}{2}(uv+vu)-av^2
 \end{gathered}
 \right).
\]
\(A_0\) は左右 Noetherian な quadratic Koszul AS 正則環で、型は
\((1,4,6,4,1)\) である。
\(g_0=uv-vu\) は左右正則正規で、
\[
 A_0/(g_0)=B_0,\qquad
 g_0u=-ug_0,\quad g_0v=-vg_0,\quad
 g_0c=cg_0,\quad g_0d=dg_0.
\]
\(\sigma\) は \(A_0\) に延長し、\(\sigma(g_0)=-g_0\) となる。
従って \(A=A_0^\sigma\) における \(g=g_0\) は中央正則二次元であり、
\[
 B_0^\sigma=A/(g).
\]
特に、指定した良い Serre 根の全軌道環に対する中央 AS4 被覆が存在する。
さらに \(E=(B_0^\sigma)^!\) の二次中心は一次元であり、
この被覆の双対線がその全体である。
\end{theorem}

\begin{proof}
全証明は相対モジュライ論文 \cite[\texttt{noncubic_central_hull}]{QuadricModuli2026} に置く。
\end{proof}

\begin{lemma}[良い noncubic 根では二次反中心が消える]
\label{lem:noncubic-no-quadratic-anticenter}
定理~\ref{thm:noncubic-central-hull} の \(E=(B_0^\sigma)^!\) について
\[
 Z(E)_2=k\chi,\qquad
 Z^-(E)_2:=\{z\in E_2\mid zx=-xz\ (x\in E_1)\}=0
\]
である。最初の式の \(\chi\) は同定 \(E=E_0^\alpha\) による像を表す。
\end{lemma}
\begin{proof}
全証明は相対モジュライ論文 \cite[\texttt{noncubic_central_hull}]{QuadricModuli2026} に置く。
\end{proof}


\begin{remark}[相対化と横断族への適用で分ける入力]
明示 noncubic 表示族上では、\(2\theta\) が単元となる任意の可換 Noetherian 基底に
上の構成を相対化できる。双対の各次数の projectivity は既存の相対 Koszul 補題による。
\(\chi\) 乗法は各幾何 fiber で単射であるから局所直和因子への単射となり、商も
任意の基底変更と両立する有限 projective 加群となる。
上の pairing の行列式は単元なので、相対 Frobenius 条件も成立する。
商に PP の相対 Koszul 定理を適用することで双対 ambient の全次数平坦性を得る。
正規元商からの Noetherian 持上げは任意の Noetherian 基底変更後にも適用できる。

横断族では、この計算だけから全軌道の平坦性や canonical heart は従わない。
ただし、被約有限型基底 \(S\) 上にその良い根の平坦な Koszul 軌道環 \(B\) があり、
\(E=B^!\) の二次・三次が階数七・八の束であると仮定する。
稠密な cubic 開部分上で同じ根の canonical hull が中央であるなら、
commutator 写像 \(E_2\to\mathcal Hom(E_1,E_3)\) の階数は稠密開で高々六である。
被約性により七次 minors は全体で零であり、上の定理は noncubic 点で階数六を与える。
従ってその点の近傍で核は任意の基底変更と両立する可逆部分束となる。
これが延ばすべき中央線であり、従来の Ore 被覆の正規線ではない。
局所環上では、次節の定理~\ref{thm:noncubic-local-central-hull} がこの入力を確認する。
全軌道平坦性、同じ根と cubic canonical 周期との比較、有限 AS4 条件を満たす
単一近傍の構成、Noetherian heart の認識は、この線の延長と別に記載する必要がある。
\end{remark}
